\chapter{Arhitektura sistema}

\section{Pregled arhitekture i tok podataka}

Sistem se sastoji od pet nezavisnih procesa: Flutter mobilne aplikacije (klijent, vozač ili dispečer), Go backend servisa (centralna tačka logike), Valhalla routing engine-a, PostgreSQL/PostGIS baze podataka i RabbitMQ posrednika poruka. Slika 4.1 prikazuje njihov međusobni odnos i osnovni tok podataka pri planiranju i praćenju jedne ture.

\begin{figure}[h]
\centering
\includegraphics[width=0.88\textwidth]{diagrams/architecture.pdf}
\caption{Arhitektura sistema — komponente i tok podataka}
\label{fig:architecture}
\end{figure}

Tok pri planiranju rute je sledeći: mobilna aplikacija šalje \texttt{POST /api/v1/routes} sa profilom vozila, polazištem i odredištem; backend mapira profil vozila u Valhalla-in \texttt{truck} costing model (poglavlje 6.1) i traži do tri alternativne rute; svaka alternativa se rangira sopstvenom funkcijom rizika (poglavlje 6.2), a za izabranu (najbolje rangiranu) rutu se, ako odstupa od geometrijski najkraćeg puta, generiše tekstualno objašnjenje (poglavlje 6.3); odgovor sa geometrijom, rizikom i objašnjenjem se vraća klijentu. Kada vozač zaista pokrene turu (\texttt{POST /api/v1/trips}, zatim prelazak u status \texttt{in\_progress}), backend objavljuje \texttt{trip.started} događaj na RabbitMQ, koji worker procès preuzima i računa predlog pauze (poglavlje 7.4); istovremeno se otvara WebSocket veza (\texttt{GET /ws/trips/\{id\}}) preko koje se, dok vozač šalje periodične GPS pozicije (\texttt{POST /api/v1/trips/\{id\}/position}), te pozicije prosleđuju svim zainteresovanim klijentima uživo — vozačevom sopstvenom ekranu i, ako postoji, njegovom dispečeru koji sistem preko sopstvene, paralelne WebSocket konekcije prati na jednoj mapi (poglavlje 7.5, 8.2).

Sistem sledi troslojnu podelu odgovornosti sličnu klasičnoj veb arhitekturi \cite{tanenbaum}: sloj prezentacije (Flutter aplikacija), sloj obrade (Go backend, uključujući scoring i explain module) i sloj perzistencije (PostgreSQL). Razlika u odnosu na uobičajenu troslojnu veb aplikaciju je što backend za svaki zahtev za rutu sam postaje klijent eksternog servisa (Valhalla), a deo obrade (predlog pauze) je pomeren u nezavisan, asinhron proces (worker), radi rasterećenja glavne HTTP niti obrade zahteva.

\section{Kontejnerizacija i orkestracija servisa}

Svi serverski delovi sistema definisani su u jednom \texttt{docker-compose.yml} fajlu i pokreću se komandom \texttt{docker compose up}. Listing 4.1 prikazuje definiciju PostgreSQL servisa, sa health-check-om koji ostali servisi (backend) koriste kao uslov za sopstveni start preko \texttt{depends\_on: condition: service\_healthy} — čime se garantuje da backend ne pokuša konekciju na bazu koja još nije spremna za prihvat konekcija.

\begin{lstlisting}[caption={Definicija PostgreSQL/PostGIS servisa u \texttt{docker-compose.yml}}]
postgres:
  image: postgis/postgis:16-3.4-alpine
  environment:
    POSTGRES_USER: hvr
    POSTGRES_PASSWORD: hvr_dev_password
    POSTGRES_DB: hvr
  ports:
    - "5433:5432"
  volumes:
    - pgdata:/var/lib/postgresql/data
  healthcheck:
    test: ["CMD-SHELL", "pg_isready -U hvr -d hvr"]
    interval: 10s
    timeout: 5s
    retries: 5
\end{lstlisting}

Tabela 4.1 sumira sve servise definisane u \texttt{docker-compose.yml}, njihovu ulogu i portove koje izlažu na host mašini.

\begin{longtable}{@{}p{0.235\textwidth}p{0.235\textwidth}p{0.235\textwidth}p{0.235\textwidth}@{}}
\caption{Servisi definisani u \texttt{docker-compose.yml}}\label{tab:4-1}\\
\toprule
\textbf{Servis} & \textbf{Slika} & \textbf{Uloga} & \textbf{Port (host)} \\
\midrule
\endhead
\texttt{postgres} & \texttt{postgis/postgis:16-3.4-alpine} & Relaciona baza + PostGIS ekstenzija & 5433 → 5432 \\
\texttt{rabbitmq} & \texttt{rabbitmq:3-management-alpine} & Posrednik poruka (AMQP) + web konzola & 5672, 15672 \\
\texttt{valhalla} & \texttt{ghcr.io/valhalla/valhalla:latest} & HTTP routing servis (\texttt{valhalla\_service}) & 8002 \\
\texttt{valhalla-build} & \texttt{ghcr.io/valhalla/valhalla:latest} & Jednokratna izgradnja grafa iz OSM ekstrakta (profil \texttt{build}, ne pokreće se pri običnom \texttt{up}) & — \\
\texttt{backend} & sopstveni Go build & REST API, WebSocket gateway, worker & 8080 \\
\bottomrule
\end{longtable}

\texttt{valhalla-build} je definisan sa Docker Compose \emph{profilom} \texttt{build}, što znači da se ne pokreće automatski pri \texttt{docker compose up} — pokreće se izričito (\texttt{docker compose --profile build up valhalla-build}) samo kada treba ponovo izgraditi graf (npr. nakon promene OSM ekstrakta, poglavlje 5), i radi sa \texttt{network\_mode: none} jer mu za samu izgradnju grafa nije potrebna mrežna konekcija, čime se smanjuje površina za eventualni bezbednosni problem tog koraka. Mobilna aplikacija i sopstveni algoritamski modul (poglavlje 6.4) nisu deo \texttt{docker-compose.yml} — Flutter aplikacija se pokreće lokalno na razvojnoj mašini ili emulatoru, a algoritamski modul je čista Go biblioteka koja se pokreće isključivo kroz testove, bez sopstvenog servisa (obrazloženo u poglavlju 6.4).

\section{Sigurnost i autentifikacija}

Autentifikacija vozača i dispečera zasniva se na JWT tokenima (RFC 7519 \cite{jwt}) potpisanim HMAC SHA-256 algoritmom (HS256). Nakon prijave (\texttt{POST /auth/login} sa korisničkim imenom i lozinkom, ili \texttt{POST /auth/google} sa Google ID tokenom), backend izdaje token koji, pored identifikatora vozača, nosi i celobrojno polje \texttt{token\_version}. Ovo polje se upoređuje sa vrednošću \texttt{token\_version} iz baze na svakom autentifikovanom zahtevu (\texttt{RequireAuth} middleware) — kada vozač zatraži odjavu sa svih uređaja (\texttt{POST /auth/logout-all}), backend samo inkrementira \texttt{token\_version} u bazi, čime svi ranije izdati tokeni (koji nose stariju vrednost) trenutno postaju nevažeći, bez potrebe za posebnom, rastućom tabelom opozvanih tokena (\emph{blocklist}) koja bi morala da se proverava na svakom zahtevu.

Lozinke se ne čuvaju u čistom tekstu, već kao bcrypt heš, algoritam koji ugrađuje nasumičnu "sol" po lozinci (čime dve identične lozinke dva različita korisnika dobijaju različit heš) i podesivi faktor rada (koji direktno određuje koliko dugo traje jedan pokušaj heširanja, pa time i koliko dugo traje jedan pokušaj brute-force napada). Nalozi kreirani preko Google prijave nemaju lozinku (\texttt{password\_hash} je \texttt{NULL}) — identifikuju se preko stabilnog Google identifikatora naloga (\texttt{google\_sub}, \texttt{sub} polje iz Google ID tokena), koji backend verifikuje samostalno preko Google-ovog javnog skupa ključeva (JWKS), bez uvođenja Firebase SDK-a kao zavisnosti.

Za obnovu zaboravljene lozinke (\texttt{POST /auth/forgot-password}) i verifikaciju email adrese (\texttt{GET /auth/verify-email}) koristi se model jednokratnog, vremenski ograničenog tokena zapisanog u posebnoj tabeli (\texttt{password\_reset\_tokens}, \texttt{email\_verification\_tokens}) sa poljima \texttt{expires\_at} i \texttt{used\_at} — link poslat emailom prestaje da važi i po isteku vremena i po jednokratnoj upotrebi, a backend razlikuje ta dva slučaja radi jasnije poruke korisniku. Slanje emaila (link za verifikaciju/reset) realizovano je preko \texttt{net/smtp} paketa iz Go standardne biblioteke, bez eksternog transakcionog email API-ja — servis je no-op (ne šalje ništa, samo zapisuje u log) ako promenljiva okruženja \texttt{SMTP\_HOST} nije podešena, čime ostatak sistema radi i bez konfigurisanog email servisa (npr. u razvojnom okruženju).

Sistem u trenutnom obliku ima namerno ograničen bezbednosni obim, obrazložen u poglavlju 10 — WebSocket \texttt{CheckOrigin} provera je trenutno propustljiva (prihvata sve origin-e), pošto je aplikacija zamišljena kao rad sa jednim poznatim, poverljivim klijentom (sopstvena Flutter aplikacija), a ne kao javni, višezakupčani (multi-tenant) servis.

\clearpage
